\cotefile{11}
\scanpage{3}
$\unsure{V|\mathcal V}$ catégorie (objets $X,Y,\ldots$) avec produits de deux \unsure{objet|objets} (\unsure{+|et} un objet final ?)

$\mathcal A\colon X\to A(X)$ foncteur cov. de $V$ dans \unsure{la|?} catégorie \deleted{\unsure{des}} \inserted{$C$} Algèbres \inserted{gr.} bivariantes, satisfaisant une condition d'échange \deleted{\unsure{entre}} \deleted{\unsure{condition}} \inserted{des $f_*$ et $g^*$} (dont le besoin exact sera \unsure{vu} ds la suite \ldots).

On va définir une nouvelle catégorie, notée $V^{\mathcal A}$, et deux foncteurs $\varphi,\psi$

\diagram{\begin{tikzcd}[ampersand replacement=\&,column sep=large]
V \arrow[r,"\varphi"] \arrow[dr,"\mathcal A"'] \& V^{\mathcal A} \arrow[r,"\psi"] \& k\text{-mod gr.}\\
\& C \arrow[ur,"\omega", "{\substack{\text{oubli}\\\text{de la structure d'\unsure{Algèbres|Algèbre|algèbre},}}}"'] \&
\end{tikzcd}}

rendant le carré écrit commutatif. De plus, le foncteur $\varphi$ induit une bijection (en fait, une identité) sur les ens d'objets, en \deleted{\unsure{d'autres termes}} \inserted{particulier}

\[
\mathrm{Ob}\,V^{\mathcal A}=\mathrm{Ob}\,V
\]

On pose.

\[
\operatorname{Hom}_{V^{\mathcal A}}(X,Y)=\mathcal A(X\times \unsure{Y|X})
\]

On définit le produit de convolution ainsi

\[
\begin{aligned}
\operatorname{Hom}_{V^{\mathcal A}}(X,Y)\times\operatorname{Hom}_{V^{\mathcal A}}(Y,Z)&\longrightarrow\operatorname{Hom}_{V^{\mathcal A}}(X,Z)\\
(u,\quad v)&\longmapsto v\circ u
\end{aligned}
\]

\[
\mathcal A(X\times Y)\times\mathcal A(Y\times Z)\longrightarrow\mathcal A(X\times Z)
\]

\diagram{\begin{tikzcd}[ampersand replacement=\&,row sep=small,column sep=small]
\&X\times Y\times Z\arrow[dl,"p_{12}"']\arrow[d,dash,"p_{13}"]\arrow[dr,"p_{23}"]\&\\
X\times Y\&X\times Z\&Y\times Z
\end{tikzcd}}

\[
\boxed{v\circ u=p_{31*}\bigl(p_{32}^{*}(\unsure{v|u})\,p_{21}^{*}(u)\bigr).}
\]

% archive sheet number: 2
